\section{余下账本事件导航：three-kingdoms}

\label{sec:remaining-event-anchors-three-kingdoms}

本节为尚未在正文设置目标的账本记录建立直接跳转。锚点显示可核日期与事件名称；其解释仍应回到账本所列来源、置信提示及既有正文的区域和材料边界。

\subsection{事件导航与来源边界}

\eventanchor{TK-LEDGER-010-5}{黄初七年五月（226年）·“曹丕去世，曹叡继位，是为魏明帝”的后续制度或军政调整}

\eventanchor{TK-LEDGER-010-6}{黄初七年五月（226年）·“曹丕去世，曹叡继位，是为魏明帝”的异源材料与影响范围的复核}

\eventanchor{TK-LEDGER-011-1}{建兴五年（227年）·“诸葛亮驻汉中并上《出师表》，以北伐准备重新组织蜀汉军政”的前期动员与资源准备}

\eventanchor{TK-LEDGER-011-2}{建兴五年（227年）·“诸葛亮驻汉中并上《出师表》，以北伐准备重新组织蜀汉军政”的命令、联盟或地方响应}

\eventanchor{TK-LEDGER-011-3}{建兴五年（227年）·诸葛亮驻汉中并上《出师表》，以北伐准备重新组织蜀汉军政}

\eventanchor{TK-LEDGER-011-4}{建兴五年（227年）·“诸葛亮驻汉中并上《出师表》，以北伐准备重新组织蜀汉军政”的执行中的空间与人员处置}

\eventanchor{TK-LEDGER-011-5}{建兴五年（227年）·“诸葛亮驻汉中并上《出师表》，以北伐准备重新组织蜀汉军政”的后续制度或军政调整}

\eventanchor{TK-LEDGER-011-6}{建兴五年（227年）·“诸葛亮驻汉中并上《出师表》，以北伐准备重新组织蜀汉军政”的异源材料与影响范围的复核}

\eventanchor{TK-LEDGER-012-1}{建兴六年春（228年）·“诸葛亮出祁山，南安、天水、安定一度响应蜀汉，魏调张郃等西援”的前期动员与资源准备}

\eventanchor{TK-LEDGER-012-2}{建兴六年春（228年）·“诸葛亮出祁山，南安、天水、安定一度响应蜀汉，魏调张郃等西援”的命令、联盟或地方响应}

\eventanchor{TK-LEDGER-012-3}{建兴六年春（228年）·诸葛亮出祁山，南安、天水、安定一度响应蜀汉，魏调张郃等西援}

\eventanchor{TK-LEDGER-012-4}{建兴六年春（228年）·“诸葛亮出祁山，南安、天水、安定一度响应蜀汉，魏调张郃等西援”的执行中的空间与人员处置}

\eventanchor{TK-LEDGER-012-5}{建兴六年春（228年）·“诸葛亮出祁山，南安、天水、安定一度响应蜀汉，魏调张郃等西援”的后续制度或军政调整}

\eventanchor{TK-LEDGER-012-6}{建兴六年春（228年）·“诸葛亮出祁山，南安、天水、安定一度响应蜀汉，魏调张郃等西援”的异源材料与影响范围的复核}

\eventanchor{TK-LEDGER-013-1}{建兴六年（228年）·“马谡失守街亭，诸葛亮撤回汉中并处置失军者”的前期动员与资源准备}

\eventanchor{TK-LEDGER-013-2}{建兴六年（228年）·“马谡失守街亭，诸葛亮撤回汉中并处置失军者”的命令、联盟或地方响应}

\eventanchor{TK-LEDGER-013-3}{建兴六年（228年）·马谡失守街亭，诸葛亮撤回汉中并处置失军者}

\eventanchor{TK-LEDGER-013-4}{建兴六年（228年）·“马谡失守街亭，诸葛亮撤回汉中并处置失军者”的执行中的空间与人员处置}

\eventanchor{TK-LEDGER-013-5}{建兴六年（228年）·“马谡失守街亭，诸葛亮撤回汉中并处置失军者”的后续制度或军政调整}

\eventanchor{TK-LEDGER-013-6}{建兴六年（228年）·“马谡失守街亭，诸葛亮撤回汉中并处置失军者”的异源材料与影响范围的复核}

\eventanchor{TK-LEDGER-014-1}{太和二年（228年）·“新城太守孟达谋反，司马懿急行军围新城，孟达为部下所杀”的前期动员与资源准备}

\eventanchor{TK-LEDGER-014-2}{太和二年（228年）·“新城太守孟达谋反，司马懿急行军围新城，孟达为部下所杀”的命令、联盟或地方响应}

\eventanchor{TK-LEDGER-014-3}{太和二年（228年）·新城太守孟达谋反，司马懿急行军围新城，孟达为部下所杀}

\eventanchor{TK-LEDGER-014-4}{太和二年（228年）·“新城太守孟达谋反，司马懿急行军围新城，孟达为部下所杀”的执行中的空间与人员处置}

\eventanchor{TK-LEDGER-014-5}{太和二年（228年）·“新城太守孟达谋反，司马懿急行军围新城，孟达为部下所杀”的后续制度或军政调整}

\eventanchor{TK-LEDGER-014-6}{太和二年（228年）·“新城太守孟达谋反，司马懿急行军围新城，孟达为部下所杀”的异源材料与影响范围的复核}

\eventanchor{TK-LEDGER-015-1}{太和二年（228年）·“陆逊在石亭击败曹休所率魏军，魏的江淮攻势受挫”的前期动员与资源准备}

\eventanchor{TK-LEDGER-015-2}{太和二年（228年）·“陆逊在石亭击败曹休所率魏军，魏的江淮攻势受挫”的命令、联盟或地方响应}

\eventanchor{TK-LEDGER-015-3}{太和二年（228年）·陆逊在石亭击败曹休所率魏军，魏的江淮攻势受挫}

\eventanchor{TK-LEDGER-015-4}{太和二年（228年）·“陆逊在石亭击败曹休所率魏军，魏的江淮攻势受挫”的执行中的空间与人员处置}

\eventanchor{TK-LEDGER-015-5}{太和二年（228年）·“陆逊在石亭击败曹休所率魏军，魏的江淮攻势受挫”的后续制度或军政调整}

\eventanchor{TK-LEDGER-015-6}{太和二年（228年）·“陆逊在石亭击败曹休所率魏军，魏的江淮攻势受挫”的异源材料与影响范围的复核}

\eventanchor{TK-LEDGER-016-1}{建兴七年（229年）·“诸葛亮取武都、阴平二郡，魏将郭淮退却，蜀汉扩展汉中西北屏障”的前期动员与资源准备}

\eventanchor{TK-LEDGER-016-2}{建兴七年（229年）·“诸葛亮取武都、阴平二郡，魏将郭淮退却，蜀汉扩展汉中西北屏障”的命令、联盟或地方响应}

\eventanchor{TK-LEDGER-016-3}{建兴七年（229年）·诸葛亮取武都、阴平二郡，魏将郭淮退却，蜀汉扩展汉中西北屏障}

\eventanchor{TK-LEDGER-016-4}{建兴七年（229年）·“诸葛亮取武都、阴平二郡，魏将郭淮退却，蜀汉扩展汉中西北屏障”的执行中的空间与人员处置}

\eventanchor{TK-LEDGER-016-5}{建兴七年（229年）·“诸葛亮取武都、阴平二郡，魏将郭淮退却，蜀汉扩展汉中西北屏障”的后续制度或军政调整}

\eventanchor{TK-LEDGER-016-6}{建兴七年（229年）·“诸葛亮取武都、阴平二郡，魏将郭淮退却，蜀汉扩展汉中西北屏障”的异源材料与影响范围的复核}

\eventanchor{TK-LEDGER-017-1}{黄龙元年四月（229年）·“孙权在武昌即皇帝位，建元黄龙，孙吴正式提出独立帝国名号”的前期动员与资源准备}

\eventanchor{TK-LEDGER-017-2}{黄龙元年四月（229年）·“孙权在武昌即皇帝位，建元黄龙，孙吴正式提出独立帝国名号”的命令、联盟或地方响应}

\eventanchor{TK-LEDGER-017-3}{黄龙元年四月（229年）·孙权在武昌即皇帝位，建元黄龙，孙吴正式提出独立帝国名号}

\eventanchor{TK-LEDGER-017-4}{黄龙元年四月（229年）·“孙权在武昌即皇帝位，建元黄龙，孙吴正式提出独立帝国名号”的执行中的空间与人员处置}

\eventanchor{TK-LEDGER-017-5}{黄龙元年四月（229年）·“孙权在武昌即皇帝位，建元黄龙，孙吴正式提出独立帝国名号”的后续制度或军政调整}

\eventanchor{TK-LEDGER-017-6}{黄龙元年四月（229年）·“孙权在武昌即皇帝位，建元黄龙，孙吴正式提出独立帝国名号”的异源材料与影响范围的复核}

\eventanchor{TK-LEDGER-018-1}{黄龙元年（229年）·“孙权将都城由武昌迁回建业，长江下游成为孙吴中枢”的前期动员与资源准备}

\eventanchor{TK-LEDGER-018-2}{黄龙元年（229年）·“孙权将都城由武昌迁回建业，长江下游成为孙吴中枢”的命令、联盟或地方响应}

\eventanchor{TK-LEDGER-018-3}{黄龙元年（229年）·孙权将都城由武昌迁回建业，长江下游成为孙吴中枢}

\eventanchor{TK-LEDGER-018-4}{黄龙元年（229年）·“孙权将都城由武昌迁回建业，长江下游成为孙吴中枢”的执行中的空间与人员处置}

\eventanchor{TK-LEDGER-018-5}{黄龙元年（229年）·“孙权将都城由武昌迁回建业，长江下游成为孙吴中枢”的后续制度或军政调整}

\eventanchor{TK-LEDGER-018-6}{黄龙元年（229年）·“孙权将都城由武昌迁回建业，长江下游成为孙吴中枢”的异源材料与影响范围的复核}

\eventanchor{TK-LEDGER-019-1}{黄龙二年（230年）·“孙权遣卫温、诸葛直率舟师求夷洲、亶洲，航海军民多病死而未获预期成果”的前期动员与资源准备}

\eventanchor{TK-LEDGER-019-2}{黄龙二年（230年）·“孙权遣卫温、诸葛直率舟师求夷洲、亶洲，航海军民多病死而未获预期成果”的命令、联盟或地方响应}

\eventanchor{TK-LEDGER-019-3}{黄龙二年（230年）·孙权遣卫温、诸葛直率舟师求夷洲、亶洲，航海军民多病死而未获预期成果}

\eventanchor{TK-LEDGER-019-4}{黄龙二年（230年）·“孙权遣卫温、诸葛直率舟师求夷洲、亶洲，航海军民多病死而未获预期成果”的执行中的空间与人员处置}

\eventanchor{TK-LEDGER-019-5}{黄龙二年（230年）·“孙权遣卫温、诸葛直率舟师求夷洲、亶洲，航海军民多病死而未获预期成果”的后续制度或军政调整}

\eventanchor{TK-LEDGER-019-6}{黄龙二年（230年）·“孙权遣卫温、诸葛直率舟师求夷洲、亶洲，航海军民多病死而未获预期成果”的异源材料与影响范围的复核}

\eventanchor{TK-LEDGER-020-1}{太和四年至五年（230年至231年）·“曹真等自关中多路准备攻蜀，栈道与山道遭大雨阻断，魏军未能深入汉中”的前期动员与资源准备}

\eventanchor{TK-LEDGER-020-2}{太和四年至五年（230年至231年）·“曹真等自关中多路准备攻蜀，栈道与山道遭大雨阻断，魏军未能深入汉中”的命令、联盟或地方响应}

\eventanchor{TK-LEDGER-020-3}{太和四年至五年（230年至231年）·曹真等自关中多路准备攻蜀，栈道与山道遭大雨阻断，魏军未能深入汉中}

\eventanchor{TK-LEDGER-020-4}{太和四年至五年（230年至231年）·“曹真等自关中多路准备攻蜀，栈道与山道遭大雨阻断，魏军未能深入汉中”的执行中的空间与人员处置}

\subsection{事件导航与来源边界}

\eventanchor{TK-LEDGER-020-5}{太和四年至五年（230年至231年）·“曹真等自关中多路准备攻蜀，栈道与山道遭大雨阻断，魏军未能深入汉中”的后续制度或军政调整}

\eventanchor{TK-LEDGER-020-6}{太和四年至五年（230年至231年）·“曹真等自关中多路准备攻蜀，栈道与山道遭大雨阻断，魏军未能深入汉中”的异源材料与影响范围的复核}

\eventanchor{TK-LEDGER-021-1}{建兴九年（231年）·“诸葛亮再次出祁山，与司马懿部相持，因粮运不继而撤军”的前期动员与资源准备}

\eventanchor{TK-LEDGER-021-2}{建兴九年（231年）·“诸葛亮再次出祁山，与司马懿部相持，因粮运不继而撤军”的命令、联盟或地方响应}

\eventanchor{TK-LEDGER-021-3}{建兴九年（231年）·诸葛亮再次出祁山，与司马懿部相持，因粮运不继而撤军}

\eventanchor{TK-LEDGER-021-4}{建兴九年（231年）·“诸葛亮再次出祁山，与司马懿部相持，因粮运不继而撤军”的执行中的空间与人员处置}

\eventanchor{TK-LEDGER-021-5}{建兴九年（231年）·“诸葛亮再次出祁山，与司马懿部相持，因粮运不继而撤军”的后续制度或军政调整}

\eventanchor{TK-LEDGER-021-6}{建兴九年（231年）·“诸葛亮再次出祁山，与司马懿部相持，因粮运不继而撤军”的异源材料与影响范围的复核}

\eventanchor{TK-LEDGER-022-1}{建兴九年（231年）·“李严因粮运失职并伪称奉诏召军，被诸葛亮弹劾废为平民”的前期动员与资源准备}

\eventanchor{TK-LEDGER-022-2}{建兴九年（231年）·“李严因粮运失职并伪称奉诏召军，被诸葛亮弹劾废为平民”的命令、联盟或地方响应}

\eventanchor{TK-LEDGER-022-3}{建兴九年（231年）·李严因粮运失职并伪称奉诏召军，被诸葛亮弹劾废为平民}

\eventanchor{TK-LEDGER-022-4}{建兴九年（231年）·“李严因粮运失职并伪称奉诏召军，被诸葛亮弹劾废为平民”的执行中的空间与人员处置}

\eventanchor{TK-LEDGER-022-5}{建兴九年（231年）·“李严因粮运失职并伪称奉诏召军，被诸葛亮弹劾废为平民”的后续制度或军政调整}

\eventanchor{TK-LEDGER-022-6}{建兴九年（231年）·“李严因粮运失职并伪称奉诏召军，被诸葛亮弹劾废为平民”的异源材料与影响范围的复核}

\eventanchor{TK-LEDGER-023-1}{太和六年（232年）·“公孙渊遣使通孙吴，孙权册拜其为燕王，试图在魏的东北边缘建立外部盟点”的前期动员与资源准备}

\eventanchor{TK-LEDGER-023-2}{太和六年（232年）·“公孙渊遣使通孙吴，孙权册拜其为燕王，试图在魏的东北边缘建立外部盟点”的命令、联盟或地方响应}

\eventanchor{TK-LEDGER-023-3}{太和六年（232年）·公孙渊遣使通孙吴，孙权册拜其为燕王，试图在魏的东北边缘建立外部盟点}

\eventanchor{TK-LEDGER-023-4}{太和六年（232年）·“公孙渊遣使通孙吴，孙权册拜其为燕王，试图在魏的东北边缘建立外部盟点”的执行中的空间与人员处置}

\eventanchor{TK-LEDGER-023-5}{太和六年（232年）·“公孙渊遣使通孙吴，孙权册拜其为燕王，试图在魏的东北边缘建立外部盟点”的后续制度或军政调整}

\eventanchor{TK-LEDGER-023-6}{太和六年（232年）·“公孙渊遣使通孙吴，孙权册拜其为燕王，试图在魏的东北边缘建立外部盟点”的异源材料与影响范围的复核}

\eventanchor{TK-LEDGER-024-1}{嘉禾二年（233年）·“公孙渊杀孙吴使者张弥、许晏并夺取所赐财物，转而向魏请和”的前期动员与资源准备}

\eventanchor{TK-LEDGER-024-2}{嘉禾二年（233年）·“公孙渊杀孙吴使者张弥、许晏并夺取所赐财物，转而向魏请和”的命令、联盟或地方响应}

\eventanchor{TK-LEDGER-024-3}{嘉禾二年（233年）·公孙渊杀孙吴使者张弥、许晏并夺取所赐财物，转而向魏请和}

\eventanchor{TK-LEDGER-024-4}{嘉禾二年（233年）·“公孙渊杀孙吴使者张弥、许晏并夺取所赐财物，转而向魏请和”的执行中的空间与人员处置}

\eventanchor{TK-LEDGER-024-5}{嘉禾二年（233年）·“公孙渊杀孙吴使者张弥、许晏并夺取所赐财物，转而向魏请和”的后续制度或军政调整}

\eventanchor{TK-LEDGER-024-6}{嘉禾二年（233年）·“公孙渊杀孙吴使者张弥、许晏并夺取所赐财物，转而向魏请和”的异源材料与影响范围的复核}

\eventanchor{TK-LEDGER-025-1}{建兴十二年（234年）·“诸葛亮出斜谷屯五丈原，与司马懿相持后病逝，蜀军撤回汉中”的前期动员与资源准备}

\eventanchor{TK-LEDGER-025-2}{建兴十二年（234年）·“诸葛亮出斜谷屯五丈原，与司马懿相持后病逝，蜀军撤回汉中”的命令、联盟或地方响应}

\eventanchor{TK-LEDGER-025-3}{建兴十二年（234年）·诸葛亮出斜谷屯五丈原，与司马懿相持后病逝，蜀军撤回汉中}

\eventanchor{TK-LEDGER-025-4}{建兴十二年（234年）·“诸葛亮出斜谷屯五丈原，与司马懿相持后病逝，蜀军撤回汉中”的执行中的空间与人员处置}

\eventanchor{TK-LEDGER-025-5}{建兴十二年（234年）·“诸葛亮出斜谷屯五丈原，与司马懿相持后病逝，蜀军撤回汉中”的后续制度或军政调整}

\eventanchor{TK-LEDGER-025-6}{建兴十二年（234年）·“诸葛亮出斜谷屯五丈原，与司马懿相持后病逝，蜀军撤回汉中”的异源材料与影响范围的复核}

\eventanchor{TK-LEDGER-026-1}{嘉禾三年（234年）·“孙权亲攻合肥新城，魏守军与援军固守，吴军撤退”的前期动员与资源准备}

\eventanchor{TK-LEDGER-026-2}{嘉禾三年（234年）·“孙权亲攻合肥新城，魏守军与援军固守，吴军撤退”的命令、联盟或地方响应}

\eventanchor{TK-LEDGER-026-3}{嘉禾三年（234年）·孙权亲攻合肥新城，魏守军与援军固守，吴军撤退}

\eventanchor{TK-LEDGER-026-4}{嘉禾三年（234年）·“孙权亲攻合肥新城，魏守军与援军固守，吴军撤退”的执行中的空间与人员处置}

\eventanchor{TK-LEDGER-026-5}{嘉禾三年（234年）·“孙权亲攻合肥新城，魏守军与援军固守，吴军撤退”的后续制度或军政调整}

\eventanchor{TK-LEDGER-026-6}{嘉禾三年（234年）·“孙权亲攻合肥新城，魏守军与援军固守，吴军撤退”的异源材料与影响范围的复核}

\eventanchor{TK-LEDGER-027-1}{建兴十二年（234年）·“诸葛亮死后，蒋琬录尚书事，继掌蜀汉中枢政务”的前期动员与资源准备}

\eventanchor{TK-LEDGER-027-2}{建兴十二年（234年）·“诸葛亮死后，蒋琬录尚书事，继掌蜀汉中枢政务”的命令、联盟或地方响应}

\eventanchor{TK-LEDGER-027-3}{建兴十二年（234年）·诸葛亮死后，蒋琬录尚书事，继掌蜀汉中枢政务}

\eventanchor{TK-LEDGER-027-4}{建兴十二年（234年）·“诸葛亮死后，蒋琬录尚书事，继掌蜀汉中枢政务”的执行中的空间与人员处置}

\eventanchor{TK-LEDGER-027-5}{建兴十二年（234年）·“诸葛亮死后，蒋琬录尚书事，继掌蜀汉中枢政务”的后续制度或军政调整}

\eventanchor{TK-LEDGER-027-6}{建兴十二年（234年）·“诸葛亮死后，蒋琬录尚书事，继掌蜀汉中枢政务”的异源材料与影响范围的复核}

\eventanchor{TK-LEDGER-028-1}{景初元年（237年）·“公孙渊拒绝入朝受命，自立燕王并与魏公开决裂”的前期动员与资源准备}

\eventanchor{TK-LEDGER-028-2}{景初元年（237年）·“公孙渊拒绝入朝受命，自立燕王并与魏公开决裂”的命令、联盟或地方响应}

\eventanchor{TK-LEDGER-028-3}{景初元年（237年）·公孙渊拒绝入朝受命，自立燕王并与魏公开决裂}

\eventanchor{TK-LEDGER-028-4}{景初元年（237年）·“公孙渊拒绝入朝受命，自立燕王并与魏公开决裂”的执行中的空间与人员处置}

\eventanchor{TK-LEDGER-028-5}{景初元年（237年）·“公孙渊拒绝入朝受命，自立燕王并与魏公开决裂”的后续制度或军政调整}

\eventanchor{TK-LEDGER-028-6}{景初元年（237年）·“公孙渊拒绝入朝受命，自立燕王并与魏公开决裂”的异源材料与影响范围的复核}
